% Part: lambda-calculus
% Chapter: syntax
% Section: substitution

\documentclass[../../../include/open-logic-section]{subfiles}

\begin{document}

\olfileid[tr]{lam}{syn}{sub}
\olsection{Yerine Koyma}

\begin{explain}
Serbest değişkenler ortam değişkenlerine yapılan göndermelerdir; dolayısıyla
serbest bir değişkenin yerine belirli bir değer kullanmak anlamlıdır. Örneğin
$\lambd[x][f x]$ içindeki $f$'yi özdeşlik fonksiyonu $\lambd[y][y]$ gibi
belirli bir terimle değiştirmek isteyebiliriz. Sonuç
$\lambd[x][(\lambd[y][y]) x]$ olur. Serbest değişkenleri lambda terimleriyle
değiştirme işlemine yerine koyma denir.
\end{explain}

\begin{defn}[Yerine koyma] \ollabel{defn:substitution}
  Bir $M$ terimindeki $x$ değişkeninin yerine $N$ terimini
  \emph{koyma} işlemi $\Subst{M}{N}{x}$, tümevarımlı olarak şöyle tanımlanır:
  \begin{enumerate}
    \item $\Subst{x}{N}{x} = N$. \ollabel{defn:substitution-1}
    \item $x \neq y$ ise $\Subst{y}{N}{x}  = y$.
      \ollabel{defn:substitution-2}
    \item $\Subst{PQ}{N}{x}  = (\Subst{P}{N}{x}) (\Subst{Q}{N}{x})$.
      \ollabel{def:substitution-3}
    \item $x \neq y$ ve $y \notin \FV{N}$ ise
      $\Subst{(\lambd[y][P])}{N}{x} =
      \lambd[y][\Subst{P}{N}{x}]$; aksi durumda tanımsızdır.
      \ollabel{defn:substitution-4}
  \end{enumerate}
\end{defn}

\begin{explain}
\olref{defn:substitution}\olref{defn:substitution-4} içinde $x\neq y$
koşulunu isteriz; çünkü $M$ içindeki $x$ değişkeninin \emph{bağlı}
geçişlerini $N$ ile değiştirmek istemeyiz. Örneğin
$\Subst{\lambd[x][x]}{y}{x}$ yerine koymasını hesaplarsak sonuç
$\lambd[x][y]$ değil, doğrudan $\lambd[x][x]$ olmalıdır.

$\lambd[y][P]$ içindeki $x$ yerine $N$ koyarken $y\notin\FV{N}$ koşulunu da
isteriz. Örneğin $\lambd[y][x]$ içinde $x$ yerine $y$ koyamayız; yani
$\Subst{\lambd[y][x]}{y}{x}$ tanımsızdır. Aksi hâlde sonuç, bir argüman alıp
onu doğrudan döndüren fonksiyonu gösteren $\lambd[y][y]$ olurdu. Oysa
$\lambd[y][x]$, her zaman $x$ terimini, daha doğrusu $x$'in gönderme yaptığı
değeri, döndüren bir fonksiyonu gösterir. Gerçekte istediğimiz sonuç, bir
argüman alıp onu kullanmadan atan ve ortam değişkeni $y$'yi döndüren bir
fonksiyondur. Bunu doğru yapmak için önce bağlı $y$ değişkenini
``yeniden adlandırmamız'' gerekir.
\end{explain}

\begin{prob}
  Aşağıdaki yerine koymaların sonucu nedir?
  \begin{enumerate}
  \item $\Subst{\lambd[y][x(\lambd[w][vwx])]}{(uv)}{x}$
  \item $\Subst{\lambd[y][x(\lambd[x][x])]}{(\lambd[y][xy])}{x}$
  \item $\Subst{y(\lambd[v][xv])}{(\lambd[y][vy])}{x}$
  \end{enumerate}
\end{prob}

\begin{thm} \ollabel{thm:notinfv}
  $x \notin \FV{M}$ ise, sol taraf tanımlı olmak koşuluyla
  $\FV{\Subst{M}{N}{x}} = \FV{M}$ olur.
\end{thm}

\begin{proof}
  $M$'nin kuruluşu üzerinde tümevarımla.
  \begin{enumerate}
  \item $M$ bir değişkendir: alıştırma.
  \item $M$, $(PQ)$ biçimindedir: alıştırma.
  \item $M$, $\lambd[y][P]$ biçiminde olsun. 
    $\Subst{\lambd[y][P]}{N}{x}$ tanımlı olduğundan
    $\lambd[y][\Subst{P}{N}{x}]$ olmak zorundadır. Dolayısıyla
    $\Subst{P}{N}{x}$ tanımlıdır; ayrıca $x\neq y$ ve
    $x\notin\FV{P}$ olur. O hâlde:
    \begin{multline*}
      \FV{\Subst{\lambd[y][P]}{N}{x}} = \\
    \begin{aligned}
      & = \FV{\lambd[y][\Subst{P}{N}{x}]} &&
      \text{\olref{defn:substitution-4} gereği}\\
      & = \FV{\Subst{P}{N}{x}} \setminus \{y\} &&
      \text{\olref[fv]{def:fv}\olref[fv]{def:fv2} gereği}\\
      & = \FV{P} \setminus \{y\} && \text{tümevarım varsayımı gereği} \\
      & = \FV{\lambd[y][P]} &&
      \text{\olref[fv]{def:fv}\olref[fv]{def:fv2} gereği.}
    \end{aligned}
    \end{multline*}
  \end{enumerate}
\end{proof}

\begin{prob}
  \olref[lam][syn][sub]{thm:notinfv} teoreminin ispatını tamamlayın.
\end{prob}

\begin{thm} \ollabel{thm:infv}
  $x \in \FV{M}$ ise, sol taraf tanımlı olmak koşuluyla
  $\FV{\Subst{M}{N}{x}} = (\FV{M} \setminus \{x\}) \cup \FV{N}$ olur.
\end{thm}

\begin{proof}
  $M$'nin kuruluşu üzerinde tümevarımla.
  \begin{enumerate}
  \item $M$ bir değişkendir: alıştırma.
  \item $M$, $PQ$ biçimindedir. $\Subst{(PQ)}{N}{x}$ tanımlı olduğundan,
    her iki yerine koyma da tanımlı olmak üzere
    $(\Subst{P}{N}{x})(\Subst{Q}{N}{x})$ olmak zorundadır. Ayrıca
    $x\in\FV{PQ}$ olduğundan $x\in\FV{P}$ ya da $x\in\FV{Q}$, belki de
    her ikisi birden, geçerlidir. Gerisi alıştırma olarak bırakılmıştır.
  \item $M$, $\lambd[y][P]$ biçiminde olsun.
    $\Subst{\lambd[y][P]}{N}{x}$ tanımlı olduğundan
    $\lambd[y][\Subst{P}{N}{x}]$ olmak zorundadır; burada
    $\Subst{P}{N}{x}$ tanımlı, $x\neq y$ ve $y\notin\FV{N}$'dir.
    $x\in\FV{\lambd[y][P]}$ olduğundan $x\in\FV{P}$ de geçerlidir. Şimdi:
    \begin{multline*}
      \FV{\Subst{(\lambd[y][P])}{N}{x}} = \\
      \begin{aligned}
      & = \FV{\lambd[y][\Subst{P}{N}{x}]} \\
      & = \FV{\Subst{P}{N}{x}} \setminus \{y\} \\
      & = ((\FV{P} \setminus \{x\}) \cup \FV{N}) \setminus \{y\}
       && \text{tümevarım varsayımı gereği} \\
      & = (\FV{P} \setminus \{x, y\}) \cup \FV{N}
       && \text{$y \notin \FV{N}$ olduğundan} \\
      & = (\FV{\lambd[y][P]} \setminus \{x\}) \cup \FV{N}.
      \end{aligned}
      \end{multline*}
  \end{enumerate}
\end{proof}

\begin{prob}
  \olref[lam][syn][sub]{thm:infv} teoreminin ispatını tamamlayın.
\end{prob}

\begin{thm}\ollabel{thm:clr}
  Sağ taraf tanımlı ve $x\notin\FV{N}$ ise
  $x \notin \FV{\Subst{M}{N}{x}}$ olur.
\end{thm}

\begin{proof}
  Alıştırma.
\end{proof}

\begin{prob}
  \olref[lam][syn][sub]{thm:clr} teoremini ispatlayın.
\end{prob}


\begin{thm}\ollabel{thm:inv}
  $\Subst{M}{y}{x}$ tanımlı ve $y \notin \FV{M}$ ise
  $\Subst{\Subst{M}{y}{x}}{x}{y} = M$ olur.
\end{thm}

\begin{proof}
  $M$'nin kuruluşu üzerinde tümevarımla.
  \begin{enumerate}
  \item $M$, bir $z$ değişkenidir: alıştırma.
  \item $M$, $(PQ)$ biçimindedir. O zaman:
    \begin{align*}
      \Subst{\Subst{(PQ)}{y}{x}}{x}{y}
      &=\Subst{((\Subst{P}{y}{x})(\Subst{Q}{y}{x}))}{x}{y} \\
      &= (\Subst{\Subst{P}{y}{x}}{x}{y})(\Subst{\Subst{Q}{y}{x}}{x}{y}) \\
      &= (PQ) \text{ tümevarım varsayımı gereği.}
    \end{align*}
  \item $M$, $\lambd[z][N]$ biçimindedir.
    $\Subst{\lambd[z][N]}{y}{x}$ tanımlı olduğundan $z\neq x$ ve
    $z\neq y$ olduğunu biliriz. Dolayısıyla:
    \begin{align*}
      \Subst{\Subst{(\lambd[z][N])}{y}{x}}{x}{y}
      & = \Subst{(\lambd[z][\Subst{N}{y}{x}])}{x}{y} \\
      & = \lambd[z][\Subst{\Subst{N}{y}{x}}{x}{y}] \\
      &= \lambd[z][N] \text{ tümevarım varsayımı gereği.}
    \end{align*}
  \end{enumerate}
\end{proof}

\begin{prob}
  \olref[lam][syn][sub]{thm:inv} teoreminin ispatını tamamlayın.
\end{prob}

\end{document}
